\documentclass[11pt]{article}
\begin{document}
\section*{Step 63 Statement}
Deflationary status: the statement below is a conditional common-refinement grounding statement, not a derivation of measured couplings or a physical parent-layer certificate.

Let REQ\_U be: ``the two child gauge-route quotients compose through one common-refinement parent closure with commuting F51 projections and compatible claim statuses.'' REQ\_U alone admits a product common-refinement control with free abelian normalization. Add the declared recognition source: ``minimal simple single-irrep common-refinement closure: among REQ_U parents, prefer a simple parent with the smallest fundamental dimension and an exact one-generation complex package.'' In the tested common-refinement parent chain this source selects the minimal simple parent class, while the larger simple parents are nonminimal extensions.

For the embedded generation, the trace computation gives
\[
\operatorname{Tr}(T_3^2)=2,\qquad
\operatorname{Tr}(Q^2)=16/3,
\]
and therefore
\[
\frac{\operatorname{Tr}(T_3^2)}{\operatorname{Tr}(Q^2)}=3/8.
\]
The normalized parent hypercharge generator has \(k_Y=5/3\), computed from \(\operatorname{Tr}(Y^2)=5/6\) in the fundamental trace convention.

Finally, the colored parent coset is the F51 refinement defect already computed in Step 47: \(\Delta_{\rm fact}(\pi_{\rm SM},\pi_U)\) has 15 finite witness pairs and the witness set is the X/Y coset used in Steps 45--46.

Exit: GROUND_landing_conditional_on_named_source. The grounding is conditional on the declared source and does not settle physical realization.
\end{document}
